% Generated by diagramas_mer.py -- Chen notation, do not edit by hand.
\begin{tikzpicture}[mer]
\node[ent] (E-CUENTAAHORRO) at (-3.00,3.00) {CUENTA\_AHORRO};
\node[entref] (E-CLIENTE) at (-3.00,9.50) {CLIENTE};
\node[entref] (E-AGENCIA) at (4.50,7.50) {AGENCIA};
\node[ent] (E-MOVIMIENTO) at (4.00,-1.50) {MOVIMIENTO};
\node[entref] (E-ASIENTOCONTABLE) at (4.00,-9.00) {ASIENTO\_CONTABLE};
\node[entref] (E-CANAL) at (-3.00,-8.50) {CANAL};
\node[ent,weak] (E-CONSUMODIARIO) at (-6.00,-2.50) {CONSUMO\_DIARIO};
\node[rel] (R-R05) at (4.00,-5.25) {CONTABILIZA};
\node[rel] (R-R06) at (0.75,5.25) {APERTURA\_EN};
\node[rel] (R-R10) at (-3.00,6.25) {ES\_TITULAR};
\node[rel] (R-R11) at (0.50,0.75) {REGISTRA};
\node[rel] (R-R12) at (0.50,-5.00) {CANALIZA};
\node[rel,ident] (R-R13) at (-4.50,0.25) {ACUMULA\_CUENTA};
\node[rel,ident] (R-R14) at (-4.50,-5.50) {ACUMULA\_CANAL};
\node[att] (A-CUENTAAHORRO-numcuenta) at (-3.89,3.96) {\underline{num\_cuenta}};
\draw[attline] (E-CUENTAAHORRO) -- (A-CUENTAAHORRO-numcuenta);
\node[att] (A-CUENTAAHORRO-cci) at (-5.02,4.36) {\underline{cci}};
\draw[attline] (E-CUENTAAHORRO) -- (A-CUENTAAHORRO-cci);
\node[att] (A-CUENTAAHORRO-moneda) at (-5.21,3.64) {moneda};
\draw[attline] (E-CUENTAAHORRO) -- (A-CUENTAAHORRO-moneda);
\node[att] (A-CUENTAAHORRO-regimen) at (-6.32,3.64) {regimen};
\draw[attline] (E-CUENTAAHORRO) -- (A-CUENTAAHORRO-regimen);
\node[att,derived] (A-CUENTAAHORRO-saldocontable) at (-5.56,2.92) {saldo\_contable};
\draw[attline] (E-CUENTAAHORRO) -- (A-CUENTAAHORRO-saldocontable);
\node[att,derived] (A-CUENTAAHORRO-saldodisponible) at (-6.86,2.20) {saldo\_disponible};
\draw[attline] (E-CUENTAAHORRO) -- (A-CUENTAAHORRO-saldodisponible);
\node[att] (A-CUENTAAHORRO-estadocuenta) at (-4.69,2.20) {estado\_cuenta};
\draw[attline] (E-CUENTAAHORRO) -- (A-CUENTAAHORRO-estadocuenta);
\node[att] (A-CUENTAAHORRO-fechaapertura) at (0.31,3.00) {fecha\_apertura};
\draw[attline] (E-CUENTAAHORRO) -- (A-CUENTAAHORRO-fechaapertura);
\node[att] (A-MOVIMIENTO-uuidtransaccion) at (1.50,-1.55) {\underline{uuid\_transaccion}};
\draw[attline] (E-MOVIMIENTO) -- (A-MOVIMIENTO-uuidtransaccion);
\node[att] (A-MOVIMIENTO-tipooperacion) at (5.56,-2.98) {tipo\_operacion};
\draw[attline] (E-MOVIMIENTO) -- (A-MOVIMIENTO-tipooperacion);
\node[att] (A-MOVIMIENTO-tipomovimiento) at (6.11,-2.20) {tipo\_movimiento};
\draw[attline] (E-MOVIMIENTO) -- (A-MOVIMIENTO-tipomovimiento);
\node[att] (A-MOVIMIENTO-monto) at (6.52,-1.48) {monto};
\draw[attline] (E-MOVIMIENTO) -- (A-MOVIMIENTO-monto);
\node[att] (A-MOVIMIENTO-glosa) at (5.61,-1.03) {glosa};
\draw[attline] (E-MOVIMIENTO) -- (A-MOVIMIENTO-glosa);
\node[att] (A-MOVIMIENTO-fechahora) at (6.21,-0.14) {fecha\_hora};
\draw[attline] (E-MOVIMIENTO) -- (A-MOVIMIENTO-fechahora);
\node[att] (A-MOVIMIENTO-estadotransaccion) at (4.10,-0.25) {estado\_transaccion};
\draw[attline] (E-MOVIMIENTO) -- (A-MOVIMIENTO-estadotransaccion);
\node[att] (A-CONSUMODIARIO-fechacorte) at (-8.28,-2.12) {\partialkey{fecha\_corte}};
\draw[attline] (E-CONSUMODIARIO) -- (A-CONSUMODIARIO-fechacorte);
\node[att,derived] (A-CONSUMODIARIO-montoacumulado) at (-3.47,-2.88) {monto\_acumulado};
\draw[attline] (E-CONSUMODIARIO) -- (A-CONSUMODIARIO-montoacumulado);
\node[att] (RA-R10-roltitular) at (-5.26,6.25) {rol\_titular};
\draw[attline] (R-R10) -- (RA-R10-roltitular);
\node[att] (RA-R10-fechavinculacion) at (-0.34,6.25) {fecha\_vinculacion};
\draw[attline] (R-R10) -- (RA-R10-fechavinculacion);
\draw[total] (E-MOVIMIENTO) -- node[card,pos=0.2] {1} (R-R05);
\draw[partial] (E-ASIENTOCONTABLE) -- node[card,pos=0.2] {1} (R-R05);
\draw[partial] (E-AGENCIA) -- node[card,pos=0.2] {1} (R-R06);
\draw[partial] (E-CUENTAAHORRO) -- node[card,pos=0.2] {N} (R-R06);
\draw[partial] (E-CLIENTE) -- node[card,pos=0.2] {N} (R-R10);
\draw[total] (E-CUENTAAHORRO) -- node[card,pos=0.2] {M} (R-R10);
\draw[partial] (E-CUENTAAHORRO) -- node[card,pos=0.2] {1} (R-R11);
\draw[total] (E-MOVIMIENTO) -- node[card,pos=0.2] {N} (R-R11);
\draw[partial] (E-CANAL) -- node[card,pos=0.2] {1} (R-R12);
\draw[partial] (E-MOVIMIENTO) -- node[card,pos=0.2] {N} (R-R12);
\draw[partial] (E-CUENTAAHORRO) -- node[card,pos=0.2] {1} (R-R13);
\draw[total] (E-CONSUMODIARIO) -- node[card,pos=0.2] {N} (R-R13);
\draw[partial] (E-CANAL) -- node[card,pos=0.2] {1} (R-R14);
\draw[total] (E-CONSUMODIARIO) -- node[card,pos=0.2] {N} (R-R14);
\end{tikzpicture}
